\documentclass[../linear-algebra.tex]{subfiles}

\begin{document}
\section{Determinants}
\subsection{10 Properties of Determinants}
For square matrices,
\begin{enumerate}
    \item $|I|=1$ for any $n\times n$ identiy matrix $I$.
    
    \item Determinant changes sign when two rows or two columns are exchanged.
    \[\begin{vmatrix} a & b \\ c & d \end{vmatrix} = - \begin{vmatrix} c & d \\ a & b \end{vmatrix} = - \begin{vmatrix} b & a \\ d & c \end{vmatrix}\]
    For permutation matrix $P$ involving $r$ row exchanges,
    \[|P| = 1 \text{ for even } r,\quad |P| = -1 \text{ for odd } r \]
    
    \item The determinant is a linear function of each row separately.\newline
    Scaling a single row of a matrix by $t$,
    \[\begin{vmatrix} ta & tb \\ c & d \end{vmatrix} = t \begin{vmatrix} a & b \\ c & d \end{vmatrix}\]
    Adding a row of one matrix to the corresponding row of another matrix,
    \[\begin{vmatrix} a+a^{\prime} & b+b^{\prime} \\ c & d \end{vmatrix} = \begin{vmatrix} a & b \\ c & d \end{vmatrix} + \begin{vmatrix} a^{\prime} & b^{\prime} \\ c & d \end{vmatrix}\]
    Note that the addition can only be distributed across a single row. The remaining rows in all the matrices involved must be unchanged.
    As a general rule,
    \[|A + B| \neq |A| + |B|\]
    
    \item If any 2 rows in matrix $A$ are equal, then $|A|=0$.
    \[\begin{vmatrix} a & b \\ a & b \\ c & d\end{vmatrix} = 0\]
    \textit{Proof.}
    Swapping the 2 identical rows in $A$, the determinant is $-|A|$ (Property 2).
    However, the matrix $A$ did not change, so $|A| = -|A| \implies |A|=0$.\newline
    Note that this applies if two rows are scalar multiples of each other.
    \[\begin{vmatrix} a & b \\ ta & tb \end{vmatrix} = t \begin{vmatrix} a & b \\ a & b \end{vmatrix} = 0\]

    \item Adding/Subtracting a multiple of one row to another row does not change $|A|$.
    \[\begin{vmatrix} a & b \\ c-la & d-lb \end{vmatrix} = \begin{vmatrix} a & b \\ c & d \end{vmatrix}\]
    for $l$ a scalar.\newline
    \textit{Proof.}
    \begin{align*}
        \begin{vmatrix} a & b \\ c-la & d-lb \end{vmatrix} &= \begin{vmatrix} a & b \\ c & d \end{vmatrix} + \begin{vmatrix} a & b \\ -la & -lb \end{vmatrix} \\
        &= \begin{vmatrix} a & b \\ c & d \end{vmatrix} - l \begin{vmatrix} a & b \\ a & b \end{vmatrix} \\
        &= \begin{vmatrix} a & b \\ c & d \end{vmatrix} - l(0)
    \end{align*}
    Note that when performing Gaussian elimination on $A$ to get $U$ (upper triangular; REF form),
    we either subtract/add multiples of one row to another (does not affect $|A|$), or swap rows (changes sign of $|A|$).
    Hence $|A| = \pm |U|$.\newline
    Note further that performing further elimination to RREF does change $|A|$,
    as we would have to scale rows (to make the pivots = 1). 

    \item If matrix $A$ has a row of zeroes, $|A|=0$.
    \[\begin{vmatrix} 0 & 0 \\ c & d \end{vmatrix} = 0\]
    \textit{Proof.}
    \[|A| = \begin{vmatrix} 0 & 0 \\ c & d \end{vmatrix} = \begin{vmatrix} 0 + c & 0 + d \\ c & d \end{vmatrix} = \begin{vmatrix} c & d \\ c & d \end{vmatrix} = 0\]

    \item If matrix $A$ is triangular, its determinant is simply the product of its main diagonal entries.
    \[|A| = a_{11}a_{22} \dots a_{nn}\]
    For purely diagonal matrix, each diagonal element can be factored out row by row until you are left with the identity matrix.
    \[\begin{vmatrix} a_{11} & 0 & 0 \\ 0 & a_{22} & 0 \\ 0 & 0 & a_{nn} \end{vmatrix} = a_{11} \begin{vmatrix} 1 & 0 & 0 \\ 0 & a_{22} & 0 \\ 0 & 0 & a_{nn} \end{vmatrix} = a_{11}a_{22} \dots a_{nn} \begin{vmatrix} 1 & 0 & 0 \\ 0 & 1 & 0 \\ 0 & 0 & 1 \end{vmatrix}\]
    Since $|I| = 1$, then we get the result above.\newline
    For upper triangular matrix $U$,
    \[|U| = \begin{vmatrix} a & b \\ 0 & d \end{vmatrix}\]
    The off-diagonal entries can be made 0 by row elimination. For example, applying $R_1 \leftarrow R_1 - \frac{b}{d}R_2$ to $U$,
    \[|U| = \begin{vmatrix} a & 0 \\ 0 & d \end{vmatrix} = ad\]
    If any diagonal element $a_{ii} = 0$, the elimination process will produce an entire row of zeros. Following Property 6, this means the determinant is $0$, and the matrix is classified as singular.

    \item $A$ is singular $\iff |A| = 0$, and $A$ is invertible $\iff |A| \neq 0$.\newline
    If $A$ is singular, the elimination process produces $U$ with a row of zeroes $\implies |U| = 0 \implies |A| = 0$.\newline
    If $A$ is invertible, $U$ will have non-zero pivots along its entire main diagonal.
    From Rule 7, the product of these pivots is non-zero, hence $|A| = \pm |U|$.\newline
    For a $2 \times 2$ matrix, and $a \neq 0$,
    \[\begin{vmatrix} a & b \\ c & d \end{vmatrix} = \begin{vmatrix} a & b \\ 0 & d - (c/a)b \end{vmatrix} = a \times \left(d - \frac{cb}{a}\right) = ad - bc\]

    \item Determinant of product $|AB| = |A||B|$.\newline
    \textit{Proof.}
    Consider the function $D(A) = \frac{|AB|}{|B|}$, if it passes these 3 rules, then $D(A)$ is a determinant.
    \begin{enumerate}
        \item Identity Matrix - If matrix $A$ is the identity matrix $I$, $D(A) = \frac{|IB|}{|B|} = \frac{|B|}{|B|} = 1$.
        \item Row exchange - If two rows of $A$ are exchanged, those exact same corresponding rows are exchanged in $AB$.
        Since exchanging rows flips the sign of a determinant, so $|AB|$ changes sign. Hence $D(A)$ changes sign too.
        \item Linearity - If one row of $A$ multiplied by a scalar $t$, that same row in $AB$ is also multiplied by $t$.
        So $|AB|$ and $D(A)$ is multiplied by $t$. The same linear behavior holds true when adding rows together.
    \end{enumerate}
    $\therefore$ $\frac{|AB|}{|B|}$ has the same properties that define $|A|$, and $\frac{|AB|}{|B|} = |A|$.\newline
    If $B$ is singular, then the product $AB$ is also singular. So $|AB| = 0$, and also $|A||B| = 0$, hence $|AB| = |A||B|$ still holds.

    \item $|A^T| = |A|$.
    For a $2 \times 2$ matrix,
    \[\begin{vmatrix} a & b \\ c & d \end{vmatrix} = \begin{vmatrix} a & c \\ b & d \end{vmatrix} = ad - bc\]
    Because of this property, all the above properties apply to columns too.
    For example,
    \begin{itemize}
        \item Swapping two columns changes sign of determinant.
        \item Adding multiple of one column to another does not affect determinant.
        \item Column of zeroes means determinant is 0.
    \end{itemize}
\end{enumerate}
\subsection{Determinants as Area/Volume}
\paragraph{Theorem} For a $2 \times 2$ matrix $A$, abs($|A|$) is the area of the parallelogram formed by the columns of $A$.\newline
For a $3 \times 3$ matrix $A$, abs($|A|$) is the volume of the parallelepiped (a 3D slanted box) formed by the columns of $A$.\newline
\textit{Proof.}\newline
Consider the $2 \times 2$ diagonal matrix,
\[A = \begin{bmatrix} a & 0 \\ 0 & d \end{bmatrix}\]
$\begin{bmatrix} a \\ 0 \end{bmatrix}$ represents a horizontal vector along the x-axis, $\begin{bmatrix} 0 \\ d \end{bmatrix}$ along the y-axis.
These vectors are perpendicular, so the shape is a perfect rectangle.
\[abs \left( \begin{vmatrix} a & 0 \\ 0 & d \end{vmatrix} \right) = abs(ad) = \text{Area of rectangle}\]
To transform any $2 \times 2$ matrix $A = [a_1\quad a_2]$ into a diagonal matrix without changing the absolute value of the determinant,
\begin{enumerate}
    \item \textbf{Swapping 2 columns} - Flips the sign of $|A|$ but absolute value is unchanged.
    \item \textbf{2D Shear, Replacing column $a_2$ with $a_2 + ca_1$} - Imagine some line $L$ passing through the origin and vector $a_1$.
    Vector $a_2$ dictates the top of the parallelogram. The line $a_2 + L$ passes through $a_2$ and is perfectly parallel to $L$.
    If we add $ca_1$ to $a_2$, the new vector $a_2 + ca_1$ simply slides along that top parallel line $a_2 + L$.
    Because the base $a_1$ hasn't changed, and the points $a_2$ and $a_2 + ca_1$ have the exact same perpendicular distance to $L$ (the height is identical), both parallelograms calculate to the exact same base $\times$ height.
    Therefore, their areas are perfectly equal and $|A|$ is unchanged.
\end{enumerate}
Similarly, for the $3\times 3$ diagonal matrix,
\[A = \begin{bmatrix} a & 0 & 0 \\ 0 & b & 0 \\ 0 & 0 & c \end{bmatrix}\]
The volume of the standard rectangular box is just $abs(|A|) = a \times b \times c$.\newline
The same logic applies to transform any $3 \times 3$ matrix $A = [a_1 \quad a_2 \quad a_3]$.
\begin{itemize}
    \item \textbf{3D Shear} - The volume of a 3D parallelepiped is calculated as the area of its base multiplied by its height.
    Let the base be the flat 2D plane formed by vectors $a_1$ and $a_3$, known mathematically as $\text{Span}\{a_1, a_3\}$.
    The height is determined by how far vector $a_2$ sticks out from that base plane. The tip of $a_2$ lives in a parallel plane defined as $a_2 + \text{Span}\{a_1, a_3\}$.
    If we alter the matrix by replacing $a_2$ with $a_2 + ca_1$, that new vector just slides around inside that top parallel plane.
    Since the vector only moves sideways within the parallel plane, any vector $a_2 + ca_1$ maintains the exact same perpendicular height as the original $a_2$.
    With the base area identical and the height unchanged, the total volume of the parallelepiped is perfectly preserved when $\begin{bmatrix} a_1 & a_2 & a_3 \end{bmatrix}$ is changed to $\begin{bmatrix} a_1 & a_2 + ca_1 & a_3 \end{bmatrix}$.
\end{itemize}
\subsection{Linear Transformations}
\paragraph{Theorem} Let $T:\mathbb{R}^2\to\mathbb{R}^2$ be a linear transformation determined by a $2 \times 2$ matrix $A$.
Then for any parallelogram $S$, the area of the transformed shape $T(S) = abs(|A|) \times$ area of $S$.\newline
\textbf{3D space}: For $T$ determined by a $3 \times 3$ matrix $A$, and $S$ a parallelepiped in $\mathbb{R}^3$,
then volume of $T(S) = abs(|A|) \times$ volume of $S$.\newline
\textit{Proof.}
First, imagine a parallelogram $S$ anchored at the origin, formed by two vectors $b_1$ and $b_2$.
We can put these vectors into their own matrix, $B = \begin{bmatrix} b_1 & b_2 \end{bmatrix}$.
Applying the transformation matrix $A$ to this shape, the new transformed shape $T(S)$ is defined by the matrix product $AB$.
The new area is $abs(|AB|) = abs(|A||B|) = abs(|A|) \times abs(|B|) = abs(|A|) \times \text{area of } S$.\newline
Then any arbitrary parallelogram can be written in the form $\mathbf{p} + S$,
where $\mathbf{p}$ is a translation vector.
\begin{align*}
    \text{area of }T(\mathbf{p} + S) &= \text{area of }(T(\mathbf{p}) + T(S)) \\
    &= \text{area of} T(S) \\
    &= abs(|A|) \times \text{area of }S \\
    &= abs(|A|) \times \text{area of }\mathbf{p} + S
\end{align*}
The above proof is similar for $3 \times 3$ case.

\vspace{0.5em}
\noindent \textbf{Expanding to Arbitrary Shapes}\newline
For example, stretching a unit circle to an ellipse.\newline
$\mathbf{u} = \begin{bmatrix} u_1 \\ u_2 \end{bmatrix}$, representing the coordinates of the points forming the original unit circle,
satisfying the circle equation $u_1^2 + u_2^2 = 1$.\newline
Diagonal matrix $A$ represents the transformation that stretches the $u_1$ axis by $a$ and the $u_2$ axis by $b$,
\[A = \begin{bmatrix} a & 0 \\ 0 & b \end{bmatrix}\]
Multiplying $\mathbf{x} = A\mathbf{u}$, then $x_1 = au_1$ and $x_2 = bu_2 \implies$ $u_1 = x_1/a$ and $u_2 = x_2/b$.
Thus we get the ellipse equation $\frac{x_1^2}{a^2} + \frac{x_2^2}{b^2} = 1$.\newline
As before, with $D$ representing the unit circle and $T$ the transformation,
\begin{align*}
    \text{area of ellipse} &= \text{area of }T(D) \\
    &= abs(|A|) \times \text{area of }D \\
    &= ab \times \pi 1^2 \\
    &= \pi ab
\end{align*}

\end{document}